\documentclass[11pt]{article}
\usepackage[a4paper,margin=25mm]{geometry}
\usepackage{amsmath,amssymb,amsthm}
\usepackage{longtable,booktabs,array}
\usepackage[hidelinks]{hyperref}
\hypersetup{pdftitle={E511781250\_1: x\textasciicircum 4+x+y\textasciicircum 3+6y+1=0 has no integer solutions}}
\newtheorem{theorem}{Theorem}
\newtheorem{lemma}[theorem]{Lemma}
% keep the space around binary operators in inline formulas
\medmuskip=4mu plus 2mu\relax
\emergencystretch=1.5em
\tolerance=400
\allowdisplaybreaks
\newcommand{\Z}{\mathbb{Z}}
\newcommand{\Q}{\mathbb{Q}}
\newcommand{\F}{\mathbb{F}}
% Legendre and Jacobi symbols: one fixed size in running text, one in displays
\newcommand{\leg}[2]{\mathchoice{\biggl(\dfrac{#1}{#2}\biggr)}{(#1/#2)}{(#1/#2)}{(#1/#2)}}
\title{\textbf{E511781250\_1}\\[4pt] {\Large The equation $x^{4}+x+y^{3}+6y+1=0$ has no integer solutions}}
\author{}
\date{}
\begin{document}
\maketitle
\vspace{-8ex}
\begin{center}
Size $h=39$, length $l=11.58$
\end{center}

\begin{theorem}\label{thm:main}
The equation
\begin{equation}\label{eq:main}
x^{4}+x+y^{3}+6y+1=0
\end{equation}
has no integer solutions.
\end{theorem}

The proof uses the cubic Hilbert reciprocity law over the field $K=\Q(\omega)$, where $\omega^2+\omega+1=0$. To a hypothetical integer solution we attach the values of 47 auxiliary polynomials $H_0,\dots,H_{46}$ with coefficients in $\Z[\omega]$, constants $c_0,\dots,c_{46}\in\Z[\omega]$, and a sum $B$ of cubic Hilbert symbols of these values, recorded by a graph with 196 weighted edges. By the reciprocity law, the local contributions $B_p$ of all primes $p$ add up to $0$ modulo $3$. We show that $B_p=0$ for $p\notin S$, compute $B_p$ for the primes $p\in S$, and obtain a contradiction. Two standard facts from class field theory are quoted without proof: the explicit formula for the cubic Hilbert symbol at the places not above $3$, and the Hilbert reciprocity law. Apart from these, the proof uses polynomial identities, which can be checked by expanding both sides, and finite computations with residues.


\section{\texorpdfstring{Cubic Hilbert symbols over $\Q(\omega)$}{Cubic Hilbert symbols over Q(w)}}\label{sec:symbols}

\paragraph{The field.} Let $K=\Q(\omega)$, where $\omega^2+\omega+1=0$. Its ring of integers is $\Z[\omega]$, and every element of $K$ can be written uniquely as $A+B\omega$ with $A,B\in\Q$. We write
\[
\langle A,B\rangle=A+B\omega,
\]
also when $A$ and $B$ are polynomials in $x$ and $y$. Then $\langle A,B\rangle\langle C,D\rangle=\langle AC-BD,\ AD+BC-BD\rangle$, and the norm is $N(\langle A,B\rangle)=A^2-AB+B^2$. In particular, a value $\langle A,B\rangle$ with $A,B\in\Z$ is zero if and only if $A=B=0$.

\paragraph{Places.} The finite places of $K$ correspond to the non-zero prime ideals of $\Z[\omega]$. A prime $p\equiv2\pmod3$ is inert: there is one place above $p$, with uniformizer $p$ and residue field $\Z[\omega]/p\Z[\omega]$ with $p^2$ elements. If $p\equiv1\pmod3$, the polynomial $t^2+t+1$ has two roots $r$ modulo $p$, and for each of them there is a place $v_r$ above $p$, at which $\omega$ is the root of $t^2+t+1$ in $\Z_p$ congruent to $r$ modulo $p$; the completion is $\Q_p$, the uniformizer is $p$, the residue field is $\F_p$, and $\omega$ reduces to $r$. The prime $3$ is ramified: $\pi=1-\omega$ generates the unique prime ideal above $3$, $\pi^2=-3\omega$, and $\Z[\omega]/\pi\Z[\omega]=\F_3$, where $\langle A,B\rangle$ reduces to $A+B$. The only infinite place of $K$ is complex. For a finite place $v$ we write $K_v$ for the completion and $v(a)$ for the normalized valuation, and we call $a$ a \emph{unit} at $v$ if $v(a)=0$.

\paragraph{Hilbert symbols.} For a place $v$ of $K$ and $a,b\in K_v^\times$, the cubic Hilbert symbol $(a,b)_v$ is a cube root of unity, see \cite[Chapter~V, \S3]{N}. We write $(a,b)_v=\omega^{\beta_v(a,b)}$ with $\beta_v(a,b)\in\Z/3\Z$, and we use the following three facts \cite[Chapter~V, \S3 and Chapter~VI, \S8]{N}.
\begin{itemize}
\item[(H1)] $\beta_v(a,bc)=\beta_v(a,b)+\beta_v(a,c)$, $\beta_v(a,b)=-\beta_v(b,a)$, and $\beta_v(a,c^3)=0$. In particular, $\beta_v(a,b)$ depends only on the classes of $a$ and $b$ in $K_v^\times/K_v^{\times3}$.
\item[(H2)] Let $v$ be a finite place not above $3$, let $q$ be the number of elements of its residue field, and let $\varpi$ be a uniformizer at $v$. Write $a=\varpi^eu$ and $b=\varpi^fw$ with units $u$ and $w$, and define $\chi(u)\in\Z/3\Z$ by $\bar u^{(q-1)/3}=\bar\omega^{\chi(u)}$, where the bar denotes the reduction to the residue field. Then $\beta_v(a,b)=e\,\chi(w)-f\,\chi(u)$. In particular, $\beta_v(a,b)=0$ if $a$ and $b$ are both units.
\item[(H3)] (Hilbert reciprocity.) For $a,b\in K^\times$, $\beta_v(a,b)=0$ for all but finitely many places $v$, and $\sum_v\beta_v(a,b)=0$, where the sum is over all places of $K$. At the complex place, $\beta_v(a,b)=0$.
\end{itemize}
The general formula in (H2) contains an additional term $ef\chi(-1)$, which vanishes because $-1=(-1)^3$ is a cube. Formula (H2) fixes the normalization of the symbols; with the opposite sign convention all values $\beta_v$ change sign, and the proof below is the same.

\paragraph{Classes at the places not above $3$.} Let $p\neq3$ be a prime and $v$ a place above $p$. For $a=p^eu\in K_v^\times$ with a unit $u$, put $c_v(a)=(e\bmod3,\ \chi(u))\in\F_3^2$. By (H2) and (H1), $c_v$ is additive, $c_v(a)$ determines the class of $a$ modulo cubes, and $\beta_v(a,b)=e\chi(w)-f\chi(u)$ for $c_v(a)=(e,\chi(u))$ and $c_v(b)=(f,\chi(w))$. We put $c_p(a)=c_v(a)\in\F_3^2$ if $p\equiv2\pmod3$, and $c_p(a)=(c_{v_r}(a),c_{v_{r'}}(a))\in\F_3^4$ if $p\equiv1\pmod3$, where $r<r'$ are the roots of $t^2+t+1$ in $\{1,\dots,p-1\}$; accordingly, $\beta_p$ is the sum of the forms $\beta_v$ over the places $v$ above $p$. For $p\equiv2\pmod3$ and a rational integer $u$ prime to $p$, we have $\bar u\in\F_p$, so that $\bar u^{(p^2-1)/3}=(\bar u^{p-1})^{(p+1)/3}=1$ and $\chi(u)=0$.

\paragraph{The place above $3$.}
\begin{lemma}\label{lem:M3}
Let $v$ be the place above $3$. Every non-zero element of $K_v$ is equal to $\pi^a\omega^b2^c(1+3\omega)^d$ times a cube, for unique $a,b,c,d\in\Z/3\Z$, and every unit congruent to $1$ modulo $9$ is a cube. Put $c_3(A)=(a,b,c,d)\in\F_3^4$. Then $c_3$ is additive, and $\beta_v(A,B)=c_3(A)^{\mathsf T}M_3\,c_3(B)$, where
\[
M_3=\begin{pmatrix}0&0&2&0\\0&0&1&2\\1&2&0&0\\0&1&0&0\end{pmatrix}\pmod3.
\]
\end{lemma}

\begin{proof}
For $\alpha\in\Z[\omega]$, the polynomial $g(s)=s+3s^2+3s^3-\alpha$ satisfies $g'(s)\equiv1\pmod\pi$, so by Hensel's lemma it has a root $s$ in the completion of $\Z[\omega]$, and then $(1+3s)^3=1+9(s+3s^2+3s^3)=1+9\alpha$. Hence every unit congruent to $1$ modulo $9$ is a cube in $K_v$, and the class of a unit modulo cubes depends only on its residue modulo $9$. There are $54$ units in $\Z[\omega]/(9)$, and a direct enumeration shows that the $27$ products $\omega^b2^c(1+3\omega)^d$ with $b,c,d\in\{0,1,2\}$ are pairwise distinct modulo $9$, also after multiplication by $-1=(-1)^3$. Hence they represent all $27$ classes of units modulo cubes. Since $\pi^3$ is a cube times a unit, the exponent of $\pi$ is the valuation modulo $3$, and the first statement follows.

By (H1), it remains to compute $\beta_v$ on the $16$ ordered pairs of the elements $\pi$, $\omega$, $2$ and $1+3\omega$. Their norms are $3$, $1$, $4$ and $7$, hence, for each such pair, all the symbols at the places not above $2$, $3$ and $7$ vanish by (H2), and, by (H3), $\beta_v$ at the place above $3$ is minus the sum of the values at the place above $2$ and at the two places above $7$. These places correspond to $\bar\omega=2$ and $\bar\omega=4$ in $\F_7$. The classes $(e,\chi)$ of the four elements at these places are
\[
\begin{array}{c|ccc}
 & 2 & \bar\omega=2 & \bar\omega=4\\\hline
\pi & (0,2) & (0,0) & (0,2)\\
\omega & (0,1) & (0,2) & (0,2)\\
2 & (1,0) & (0,2) & (0,1)\\
1+3\omega & (0,2) & (1,0) & (0,0)
\end{array}
\]
where, at the places above $7$, $\chi(u)$ is defined by $\bar u^2=\bar\omega^{\chi(u)}$ in $\F_7$, and $7$ is a uniformizer. Formula (H2) and the sum over these three places give $M_3$.
\end{proof}

We write $\beta_3$ for the form of Lemma~\ref{lem:M3}. Since $3=-\omega^2\pi^2$ and $4=2^2$, we have $c_3(3)=(2,2,0,0)$ and $c_3(4)=(0,0,2,0)$.

\paragraph{Residue classes.} The local computations use the following description of the classes of the values of a polynomial on a residue class.

\begin{lemma}\label{lem:residue}
Let $p$ be a prime, $k\geq1$, and let $z\in\Z[\omega]$ be non-zero with $z\equiv\zeta\pmod{p^k}$, where $\zeta=\langle A,B\rangle$ with $0\leq A,B<p^k$. Then $c_p(z)\in b+D$ for the vector $b$ and the subspace $D$ given as follows.
\begin{itemize}
\item[(i)] $p\equiv2\pmod3$. If $\zeta\neq0$, then $b=c_p(\zeta)$ and $D=0$. If $\zeta=0$, then $b=0$, and $D$ is $\F_3^2$, or the span of $(1,0)$ if $z\in\Z$.
\item[(ii)] $p\equiv1\pmod3$. For each root $r$, let $\rho$ be the root of $t^2+t+1$ modulo $p^k$ congruent to $r$. If $A+B\rho\not\equiv0\pmod{p^k}$, the two coordinates of $c_{v_r}(z)$ are those of $c_{v_r}(\zeta)$; otherwise these two coordinates are arbitrary.
\item[(iii)] $p=3$, $z\notin\Z$. If $\zeta\neq0$, let $a=v(\zeta)<2k$ be its valuation at $\pi$ and $m=2k-a$. Then $b=c_3(\zeta)$, and $D=0$ if $m\geq4$, while $D=D_m$ if $m\leq3$, where $D_1$ is spanned by $(0,1,0,0)$, $(0,0,1,0)$ and $(0,0,0,1)$, $D_2$ by $(0,0,1,0)$ and $(0,0,0,1)$, and $D_3$ by $(0,0,1,1)$. If $\zeta=0$, then $b=0$ and $D=\F_3^4$.
\item[(iv)] $p=3$, $z\in\Z$, so that $B=0$. If $3^k\nmid A$, let $j=v_3(A)$. Then $b=c_3(A)$ and $D=0$ if $k-j\geq2$, while $b=j\,c_3(3)$ and $D$ is spanned by $c_3(4)$ if $k-j=1$. If $3^k\mid A$, then $b=0$ and $D$ is spanned by $c_3(3)$ and $c_3(4)$.
\end{itemize}
\end{lemma}

\begin{proof}
(i) If $\zeta\neq0$, then $e=v(z)=v(\zeta)<k$ and $z/p^e\equiv\zeta/p^e\pmod p$, so that $\chi$ is determined. If $z\in\Z$, its unit part is a rational integer, and $\chi=0$. (ii) At $v_r$, $z$ and $\zeta$ are the $p$-adic numbers $A'+B'\omega$ and $A+B\omega$ with $\omega\equiv\rho\pmod{p^k}$, and the argument of (i) applies. (iii) If $\zeta\neq0$, then $v(z)=a$, and $z/\pi^a\equiv\zeta/\pi^a\pmod{\pi^m}$. By Lemma~\ref{lem:M3}, the class of a unit depends only on its residue modulo $9=\omega\pi^4$; the units congruent to $1$ modulo $\pi^m$ form a subgroup, and a direct enumeration of their residues modulo $9$ shows that their classes form $D_m$. (iv) If $3^k\nmid A$, then $z=3^ju$ with a rational unit $u\equiv A/3^j\pmod{3^{k-j}}$. The class of a rational unit depends only on its residue modulo $9$, the units $\pm1$ are cubes, and the rational units congruent to $1$ modulo $3$ are $4^s$ times a unit congruent to $1$ modulo $9$; hence the classes are as stated, and if $3^k\mid A$ the valuation $j$ is also unknown.
\end{proof}

\begin{lemma}\label{lem:envelope}
Let $p$ be a prime, and consider vertices $i$ with vectors $b_i$, subspaces $D_i$ and vectors $\gamma_i$, and a set of edges $(i,j)$ with $i<j$ and weights $w_{ij}\in\{1,2\}$. Put $\rho_{ij}=w_{ij}$ and $\rho_{ji}=-w_{ij}$ for the edges, and $\sigma_i=\gamma_i+\sum_j\rho_{ij}b_j$. Suppose that $\beta_p(d,\sigma_i)=0$ for every $i$ and every $d\in D_i$, and that $\beta_p(d,d')=0$ for every edge $(i,j)$ and all $d\in D_i$, $d'\in D_j$. Then
\[
\sum_{(i,j)}w_{ij}\beta_p(c_i,c_j)+\sum_i\beta_p(c_i,\gamma_i)=\sum_{(i,j)}w_{ij}\beta_p(b_i,b_j)+\sum_i\beta_p(b_i,\gamma_i)
\]
whenever $c_i\in b_i+D_i$ for every $i$. By bilinearity, it suffices to check the hypotheses for $d$ and $d'$ in spanning sets of $D_i$ and $D_j$.
\end{lemma}

\begin{proof}
Write $c_i=b_i+d_i$ with $d_i\in D_i$. As $\beta_p$ is bilinear and $\beta_p(u,v)=-\beta_p(v,u)$, the left side equals the right side plus $\sum_i\beta_p(d_i,\sigma_i)+\sum_{(i,j)}w_{ij}\beta_p(d_i,d_j)$, and both sums vanish.
\end{proof}
\section{The graph}

Put $F(x,y)=x^{4}+x+y^{3}+6y+1$, so that \eqref{eq:main} is the equation $F(x,y)=0$. The certificate consists of 47 polynomials $H_i$ with coefficients in $\Z[\omega]$ and constants $c_i\in\Z[\omega]$, listed in Table~\ref{tab:vert1}, and of 196 \emph{edges} $(i,j)$ with $i<j$ and weights $w_{ij}\in\{1,2\}$, listed below. For an edge $(i,j)$ we put $\rho_{ij}=w_{ij}$ and $\rho_{ji}=-w_{ij}\bmod 3$, and $\rho_{ij}=0$ if $i$ and $j$ are not joined by an edge.
The edges $(i,j;w_{ij})$ are $(0,7;2)$, $(0,8;2)$, $(0,11;1)$, $(0,13;1)$, $(0,15;2)$, $(0,16;1)$, $(0,20;2)$, $(0,23;1)$, $(0,25;1)$, $(0,29;1)$, $(0,40;2)$, $(0,43;1)$, $(0,45;1)$, $(1,5;1)$, $(1,7;1)$, $(1,11;1)$, $(1,16;1)$, $(1,24;1)$, $(1,27;2)$, $(1,28;1)$, $(1,33;1)$, $(1,43;1)$, $(2,8;2)$, $(2,15;1)$, $(2,20;2)$, $(2,22;1)$, $(2,24;2)$, $(2,27;1)$, $(2,29;1)$, $(2,30;1)$, $(2,33;2)$, $(2,40;1)$, $(2,41;2)$, $(2,42;2)$, $(2,43;1)$, $(3,7;2)$, $(3,14;2)$, $(3,15;1)$, $(3,29;1)$, $(3,30;1)$, $(3,36;1)$, $(3,39;2)$, $(3,40;1)$, $(3,41;2)$, $(3,44;1)$, $(4,7;2)$, $(4,10;1)$, $(4,11;1)$, $(4,14;2)$, $(4,27;2)$, $(4,35;2)$, $(4,41;1)$, $(5,14;1)$, $(5,32;2)$, $(5,34;2)$, $(5,35;2)$, $(6,8;1)$, $(6,18;2)$, $(6,26;2)$, $(6,37;2)$, $(6,42;2)$, $(6,43;1)$, $(7,9;2)$, $(7,17;1)$, $(7,18;1)$, $(7,19;1)$, $(7,21;2)$, $(7,24;1)$, $(7,25;1)$, $(7,31;1)$, $(7,37;1)$, $(7,38;1)$, $(8,9;2)$, $(8,13;1)$, $(8,18;1)$, $(8,24;2)$, $(8,26;2)$, $(8,34;1)$, $(8,39;1)$, $(8,40;1)$, $(8,46;2)$, $(9,25;1)$, $(9,41;1)$, $(10,19;2)$, $(10,28;1)$, $(10,29;2)$, $(10,35;1)$, $(10,39;2)$, $(10,40;2)$, $(10,42;1)$, $(10,44;2)$, $(11,18;1)$, $(11,24;2)$, $(11,28;1)$, $(11,35;1)$, $(11,43;2)$, $(12,22;2)$, $(12,25;2)$, $(12,29;2)$, $(12,38;1)$, $(12,44;1)$, $(13,17;1)$, $(13,19;1)$, $(13,24;1)$, $(13,26;2)$, $(13,32;2)$, $(13,34;2)$, $(13,35;2)$, $(13,36;2)$, $(13,39;1)$, $(13,41;1)$, $(13,46;2)$, $(14,17;1)$, $(14,20;1)$, $(14,22;1)$, $(14,27;2)$, $(14,29;2)$, $(14,33;2)$, $(14,46;2)$, $(15,22;2)$, $(15,37;1)$, $(15,39;2)$, $(15,41;1)$, $(16,23;2)$, $(16,32;2)$, $(16,33;2)$, $(16,35;2)$, $(16,38;1)$, $(16,42;2)$, $(17,21;2)$, $(17,22;1)$, $(17,29;2)$, $(17,31;1)$, $(17,32;2)$, $(17,36;2)$, $(18,31;1)$, $(18,37;1)$, $(19,32;1)$, $(19,42;2)$, $(20,25;2)$, $(20,26;1)$, $(20,29;2)$, $(20,30;2)$, $(20,32;2)$, $(20,36;2)$, $(21,25;1)$, $(21,26;1)$, $(21,28;2)$, $(21,30;2)$, $(21,35;2)$, $(21,38;1)$, $(21,43;2)$, $(21,44;1)$, $(21,45;1)$, $(22,41;2)$, $(23,31;1)$, $(23,32;2)$, $(23,37;2)$, $(23,40;2)$, $(24,25;2)$, $(24,35;2)$, $(24,36;1)$, $(24,39;2)$, $(24,44;1)$, $(25,38;1)$, $(25,44;2)$, $(26,29;1)$, $(27,45;2)$, $(28,43;2)$, $(29,36;1)$, $(29,38;2)$, $(29,44;2)$, $(30,32;1)$, $(30,36;1)$, $(30,37;2)$, $(31,40;2)$, $(32,35;1)$, $(32,41;1)$, $(32,43;1)$, $(32,45;1)$, $(32,46;1)$, $(33,35;2)$, $(33,38;2)$, $(33,40;1)$, $(33,43;2)$, $(34,41;2)$, $(34,43;2)$, $(35,39;2)$, $(35,41;2)$, $(35,42;2)$, $(36,37;1)$, $(38,40;1)$, $(39,46;1)$, $(40,45;2)$, $(42,45;1)$, $(43,45;1)$.
Let $(\xi,\eta)$ denote $(x,y)$ or $(y,x)$. We call a vertex $i$ \emph{linear} if $H_i=\varepsilon_i(\eta-r_i(\xi))$ for a unit $\varepsilon_i$ of $\Z[\omega]$ and a polynomial $r_i$ with coefficients in $\Z[\omega]$; the last column of Table~\ref{tab:vert1} gives $\eta=r_i(\xi)$ for the linear vertices. For a linear vertex $i$, we put $f_i(t)=F$ and $g_{ij}(t)=H_j$ evaluated at $\xi=t$, $\eta=r_i(t)$; these are polynomials in $t$ with coefficients in $\Z[\omega]$.
For every linear vertex $i$, Table~\ref{tab:star} gives an integer $D_i\geq1$ and a polynomial $w_i(t)$ with coefficients in $\Z[\omega]$ such that
\begin{equation}\label{eq:starlin}
D_i^3\,c_i\prod_jg_{ij}^{\rho_{ij}}-w_i^3=f_iQ_i
\end{equation}
for a polynomial $Q_i(t)$ with coefficients in $\Z[\omega]$, where $\rho_{ij}\in\{0,1,2\}$; $Q_i$ is the quotient of the division of the left side by $f_i$, and the remainder is zero.
For every edge $(i,j)$ with a linear endpoint, Table~\ref{tab:sep} gives the first linear endpoint $k\in\{i,j\}$ and the factorization of the norm of the resultant $R_{ij}=\operatorname{Res}_t(f_k,g_{kl})$, where $\{k,l\}=\{i,j\}$; it is non-zero.
These identities can be checked by expanding both sides and by polynomial division, using $\omega^2=-1-\omega$.

\section{Proof of Theorem~\ref{thm:main}}

Suppose that $(x,y)\in\Z^2$ is a solution of \eqref{eq:main}, so that $F(x,y)=0$; from now on the polynomials denote their values at $(x,y)$, which lie in $\Z[\omega]$.
\paragraph{Step 1: the values are non-zero.} For $i\in\{0,\allowbreak 1,\allowbreak 2,\allowbreak 4,\allowbreak 5,\allowbreak 6,\allowbreak 11,\allowbreak 12,\allowbreak 14,\allowbreak 15,\allowbreak 16,\allowbreak 17,\allowbreak 18,\allowbreak 21,\allowbreak 22,\allowbreak 23,\allowbreak 27,\allowbreak 31,\allowbreak 33,\allowbreak 37,\allowbreak 42\}$, there is no residue pair $(a,b)$ modulo $2$ with $F(a,b)\equiv0$ and $H_i(a,b)\equiv0\pmod{2}$, as one checks for the $2^2$ residue pairs; hence $H_i\neq0$. For $i\in\{20,\allowbreak 25,\allowbreak 26,\allowbreak 32,\allowbreak 34,\allowbreak 38,\allowbreak 39,\allowbreak 40,\allowbreak 41,\allowbreak 44,\allowbreak 45\}$, there is no residue pair $(a,b)$ modulo $3$ with $F(a,b)\equiv0$ and $H_i(a,b)\equiv0\pmod{3}$, as one checks for the $3^2$ residue pairs; hence $H_i\neq0$. For $i\in\{8,\allowbreak 19,\allowbreak 28,\allowbreak 35,\allowbreak 36,\allowbreak 43,\allowbreak 46\}$, there is no residue pair $(a,b)$ modulo $4$ with $F(a,b)\equiv0$ and $H_i(a,b)\equiv0\pmod{4}$, as one checks for the $4^2$ residue pairs; hence $H_i\neq0$. For $i\in\{3,\allowbreak 9,\allowbreak 10,\allowbreak 29,\allowbreak 30\}$, there is no residue pair $(a,b)$ modulo $5$ with $F(a,b)\equiv0$ and $H_i(a,b)\equiv0\pmod{5}$, as one checks for the $5^2$ residue pairs; hence $H_i\neq0$. For $i\in\{7,\allowbreak 13\}$, there is no residue pair $(a,b)$ modulo $9$ with $F(a,b)\equiv0$ and $H_i(a,b)\equiv0\pmod{9}$, as one checks for the $9^2$ residue pairs; hence $H_i\neq0$. For $i\in\{24\}$, there is no residue pair $(a,b)$ modulo $11$ with $F(a,b)\equiv0$ and $H_i(a,b)\equiv0\pmod{11}$, as one checks for the $11^2$ residue pairs; hence $H_i\neq0$. Also, all $c_i$ are non-zero.

\paragraph{Step 2: reciprocity.} For every place $v$ of $K$, put
\begin{equation}\label{eq:B}
B_v=\sum_{(i,j)}w_{ij}\,\beta_v(H_i,H_j)+\sum_i\beta_v(H_i,c_i)\in\Z/3\Z,
\end{equation}
the first sum over the edges, and, for a rational prime $p$, put $B_p=\sum_{v\mid p}B_v$. Applying (H3) to every pair $(H_i,H_j)$ with an edge, multiplied by $w_{ij}$, and to every pair $(H_i,c_i)$, and adding, we find that $B_v=0$ for all but finitely many $v$ and $\sum_pB_p=0$, as the complex place contributes $0$. By (H1) and the definition of $c_p$, $B_p=\sum_{(i,j)}w_{ij}\beta_p(c_p(H_i),c_p(H_j))+\sum_i\beta_p(c_p(H_i),c_p(c_i))$, where $\beta_p$ is the form of Lemma~\ref{lem:M3} if $p=3$.

\paragraph{Step 3: the primes outside $S$.} Let $S=\{2,3,7,13\}$. It contains $3$ and the primes dividing the norms of the constants $c_i$, of the $D_i$, of the $R_{ij}$. Let $p\notin S$, and let $v$ be a place above $p$, with residue field $k_v$; then all these numbers are units at $v$. First, no edge $(i,j)$ has both values divisible by $v$. Indeed, suppose that $v$ divides $H_i$ and $H_j$. If $R_{ij}$ is defined, let $k$ be the linear endpoint used for it and $\{k,l\}=\{i,j\}$, and put $t=\xi$. As $H_k=\varepsilon_k(\eta-r_k(\xi))$, we have $\eta\equiv r_k(t)$ modulo $v$, hence $f_k(t)\equiv F(x,y)=0$ and $g_{kl}(t)\equiv H_l\equiv0$ modulo $v$. The resultant can be written as $R_{ij}=Af_k+Cg_{kl}$ with polynomials $A$ and $C$ with coefficients in $\Z[\omega]$, so $v$ divides $R_{ij}$, a contradiction. By (H2), a term of \eqref{eq:B} in which all arguments are units vanishes. Now let $v$ divide $H_i$. Collecting the terms of \eqref{eq:B} which contain $H_i$, and using (H1) and $\rho_{ji}=-w_{ij}$ for the edges $(j,i)$ with $j<i$, we obtain $\beta_v(H_i,\Pi_i)$ with $\Pi_i=c_i\prod_jH_j^{\rho_{ij}}$, which is a unit. We claim that $D_i^3\Pi_i\equiv W^3$ modulo $v$ for some $W\in\Z[\omega]$. If $i$ is linear, put $t=\xi$; then $g_{ij}(t)\equiv H_j$ and $f_i(t)\equiv0$ modulo $v$ as above, and \eqref{eq:starlin} gives the claim with $W=w_i(t)$. Hence the reduction of $\Pi_i$ is a non-zero cube in $k_v$, $\chi(\Pi_i)=0$, and $\beta_v(H_i,\Pi_i)=v(H_i)\chi(\Pi_i)=0$ by (H2). Hence $B_v=0$, and $B_p=0$.

\paragraph{Step 4: the primes in $S$.} Let $p\in S$. Here $2$ is inert, the places above $7$ are $v_{2}$ and $v_{4}$, and the places above $13$ are $v_{3}$ and $v_{9}$. Suppose that $x\equiv a$ and $y\equiv b\pmod{p^k}$. Then $H_i\equiv H_i(a,b)\pmod{p^k}$, and Lemma~\ref{lem:residue}, applied with $\zeta$ the residue of $H_i(a,b)$, gives a vector $b_i$ and a subspace $D_i$ with $c_p(H_i)\in b_i+D_i$; when $H_i$ has rational coefficients, $H_i\in\Z$, and we use this in (i) and (iv). With $\gamma_i=c_p(c_i)$, if the hypotheses of Lemma~\ref{lem:envelope} hold, then $B_p$ is determined by $a$, $b$ and $k$. Starting from the solutions of \eqref{eq:main} modulo $p$, we refine a residue class $x\equiv a$, $y\equiv b\pmod{p^k}$ to the classes modulo $p^{k+1}$ that contain solutions of \eqref{eq:main} modulo $p^{k+1}$, until these hypotheses hold, or until no solution remains. Table~\ref{tab:local} lists the resulting terminal classes and the value of $B_p$ in each; every integer solution lies in one of them. The result is $B_{2}=0$, $B_{3}=0$, $B_{7}=1$, $B_{13}=0$.

\paragraph{Step 5: contradiction.} By Steps~2 and~3, $\sum_{p\in S}B_p=0$. But by Step~4, $\sum_{p\in S}B_p=1$ in $\Z/3\Z$. This contradiction proves the theorem.
\qed


\begin{thebibliography}{9}
\bibitem{N} J.~Neukirch, \emph{Algebraic Number Theory}, Grundlehren der mathematischen Wissenschaften 322, Springer, Berlin, 1999.
\end{thebibliography}

\clearpage
\begingroup\small
\setlength{\LTcapwidth}{\textwidth}
\begin{longtable}{@{}rlll@{}}
\caption{The polynomials $H_i$, the constants $c_i$ and, for the linear vertices, $\eta=r_i(\xi)$. Here $\langle A,B\rangle=A+B\omega$.}\label{tab:vert1}\\
\hline $i$ & $H_i$ & $c_i$ & chart\\\hline\endfirsthead\hline $i$ & $H_i$ & $c_i$ & chart\\\hline\endhead\hline\endfoot
0 & $\langle -y+1,1\rangle$ & $\langle -1,-1\rangle$ & $y=\langle 1,1\rangle$\\
1 & $\langle -y,1\rangle$ & $\langle -4,-4\rangle$ & $y=\langle 0,1\rangle$\\
2 & $\langle -y,-1\rangle$ & $\langle 0,-6\rangle$ & $y=\langle 0,-1\rangle$\\
3 & $\langle -y-1,0\rangle$ & $\langle 1,0\rangle$ & $y=\langle -1,0\rangle$\\
4 & $\langle -y-1,-1\rangle$ & $\langle -4,-4\rangle$ & $y=\langle -1,-1\rangle$\\
5 & $\langle -2x-y,-x\rangle$ & $\langle 1,0\rangle$ & $y=\langle -2x,-x\rangle$\\
6 & $\langle -x-y+1,x+1\rangle$ & $\langle -4,8\rangle$ & $y=\langle -x+1,x+1\rangle$\\
7 & $\langle x^{2}-y+1,0\rangle$ & $\langle -3,0\rangle$ & $y=\langle x^{2}+1,0\rangle$\\
8 & $\langle -x^{2}-x-y-1,-x^{2}-x\rangle$ & $\langle -1,-1\rangle$ & $y=\langle -x^{2}-x-1,-x^{2}-x\rangle$\\
9 & $\langle x,0\rangle$ & $\langle -4,-4\rangle$ & $x=\langle 0,0\rangle$\\
10 & $\langle x+1,0\rangle$ & $\langle 1,-1\rangle$ & $x=\langle -1,0\rangle$\\
11 & $\langle x,1\rangle$ & $\langle 4,-4\rangle$ & $x=\langle 0,-1\rangle$\\
12 & $\langle x,-1\rangle$ & $\langle 0,-2\rangle$ & $x=\langle 0,1\rangle$\\
13 & $\langle x-1,0\rangle$ & $\langle -2,-4\rangle$ & $x=\langle 1,0\rangle$\\
14 & $\langle x-1,-1\rangle$ & $\langle 0,-3\rangle$ & $x=\langle 1,1\rangle$\\
15 & $\langle x-1,1\rangle$ & $\langle -12,0\rangle$ & $x=\langle 1,-1\rangle$\\
16 & $\langle x-2,-1\rangle$ & $\langle 14,30\rangle$ & $x=\langle 2,1\rangle$\\
17 & $\langle x,-y\rangle$ & $\langle 4,0\rangle$ & $y=\langle -x,-x\rangle$\\
18 & $\langle x+y+1,1\rangle$ & $\langle 0,-12\rangle$ & $y=\langle -x-1,-1\rangle$\\
19 & $\langle x+y+1,y+1\rangle$ & $\langle 1,-1\rangle$ & $y=\langle -1,x\rangle$\\
20 & $\langle x-y+1,0\rangle$ & $\langle 2,0\rangle$ & $y=\langle x+1,0\rangle$\\
21 & $\langle x+y,-1\rangle$ & $\langle -4,-4\rangle$ & $y=\langle -x,1\rangle$\\
22 & $\langle x-y-1,-y\rangle$ & $\langle 4,0\rangle$ & $y=\langle 0,-x+1\rangle$\\
23 & $\langle x-1,y\rangle$ & $\langle 4,-4\rangle$ & $y=\langle x-1,x-1\rangle$\\
24 & $\langle x+y-1,0\rangle$ & $\langle 0,-4\rangle$ & $y=\langle -x+1,0\rangle$\\
25 & $\langle x-1,y-1\rangle$ & $\langle 1,0\rangle$ & $y=\langle x,x-1\rangle$\\
26 & $\langle x-y+1,-y-1\rangle$ & $\langle -5,3\rangle$ & $y=\langle -1,-x-2\rangle$\\
27 & $\langle x-y+1,-1\rangle$ & $\langle -3,0\rangle$ & $y=\langle x+1,-1\rangle$\\
28 & $\langle x+1,-y-1\rangle$ & $\langle -8,-4\rangle$ & $y=\langle -x-2,-x-1\rangle$\\
29 & $\langle x+y-1,y+1\rangle$ & $\langle 2,0\rangle$ & $y=\langle -1,x-2\rangle$\\
30 & $\langle x-2,-y-1\rangle$ & $\langle 1,0\rangle$ & $y=\langle -x+1,-x+2\rangle$\\
31 & $\langle x+1,1\rangle$ & $\langle 4,0\rangle$ & $x=\langle -1,-1\rangle$\\
32 & $\langle -y+1,-x\rangle$ & $\langle 4,0\rangle$ & $y=\langle 1,-x\rangle$\\
33 & $\langle x-y,-y\rangle$ & $\langle 0,-4\rangle$ & $y=\langle 0,-x\rangle$\\
34 & $\langle -x-y,0\rangle$ & $\langle 2,0\rangle$ & $y=\langle -x,0\rangle$\\
35 & $\langle -x-y-1,-x^{2}\rangle$ & $\langle -1,-2\rangle$ & $y=\langle -x-1,-x^{2}\rangle$\\
36 & $\langle -x^{2}-x-y-1,0\rangle$ & $\langle -6,-6\rangle$ & $y=\langle -x^{2}-x-1,0\rangle$\\
37 & $\langle -y,x\rangle$ & $\langle 1,-1\rangle$ & $y=\langle 0,x\rangle$\\
38 & $\langle -y-1,-2\rangle$ & $\langle -1,-1\rangle$ & $y=\langle -1,-2\rangle$\\
39 & $\langle x-y+1,x\rangle$ & $\langle 0,-12\rangle$ & $y=\langle x+1,x\rangle$\\
40 & $\langle x+1,-y+1\rangle$ & $\langle 1,0\rangle$ & $y=\langle -x,-x-1\rangle$\\
41 & $\langle -y+1,x+1\rangle$ & $\langle 1,0\rangle$ & $y=\langle 1,x+1\rangle$\\
42 & $\langle x-y+1,x+1\rangle$ & $\langle -2,-2\rangle$ & $y=\langle x+1,x+1\rangle$\\
43 & $\langle -y-1,x^{2}+x\rangle$ & $\langle 0,-1\rangle$ & $y=\langle -1,x^{2}+x\rangle$\\
44 & $\langle -x-y,-x+1\rangle$ & $\langle 1,0\rangle$ & $y=\langle -x,-x+1\rangle$\\
45 & $\langle x+2,2\rangle$ & $\langle 4,16\rangle$ & $x=\langle -2,-2\rangle$\\
46 & $\langle x+y+2,0\rangle$ & $\langle 0,-4\rangle$ & $y=\langle -x-2,0\rangle$\\
\end{longtable}
\endgroup

\begingroup\scriptsize
\setlength{\LTcapwidth}{\textwidth}
\begin{longtable}{@{}ll@{}}
\caption{The data of the identities \eqref{eq:starlin}; the entry $j{:}e$ means $\rho_{ij}=e\neq0$.}\label{tab:star}\\
\hline\endfirsthead\hline\endhead\hline\endfoot
\multicolumn{2}{@{}p{0.97\textwidth}@{}}{$i=0$, $D_i=1$, $\rho_{ij}$: $7{:}2, \allowbreak 8{:}2, \allowbreak 11{:}1, \allowbreak 13{:}1, \allowbreak 15{:}2, \allowbreak 16{:}1, \allowbreak 20{:}2, \allowbreak 23{:}1, \allowbreak 25{:}1, \allowbreak 29{:}1, \allowbreak 40{:}2, \allowbreak 43{:}1, \allowbreak 45{:}1$}\\
$w_i$ & $\langle 14t^{3}-8t^{2}+8t-54,$\\
 & $\quad 10t^{3}+2t^{2}+10t\rangle$\\
\addlinespace[2pt]
\multicolumn{2}{@{}p{0.97\textwidth}@{}}{$i=1$, $D_i=1$, $\rho_{ij}$: $5{:}1, \allowbreak 7{:}1, \allowbreak 11{:}1, \allowbreak 16{:}1, \allowbreak 24{:}1, \allowbreak 27{:}2, \allowbreak 28{:}1, \allowbreak 33{:}1, \allowbreak 43{:}1$}\\
$w_i$ & $\langle -2t^{3}-4t^{2}-6t-4,$\\
 & $\quad 2t^{2}-4t-6\rangle$\\
\addlinespace[2pt]
\multicolumn{2}{@{}p{0.97\textwidth}@{}}{$i=2$, $D_i=1$, $\rho_{ij}$: $8{:}2, \allowbreak 15{:}1, \allowbreak 20{:}2, \allowbreak 22{:}1, \allowbreak 24{:}2, \allowbreak 27{:}1, \allowbreak 29{:}1, \allowbreak 30{:}1, \allowbreak 33{:}2, \allowbreak 40{:}1, \allowbreak 41{:}2, \allowbreak 42{:}2, \allowbreak 43{:}1$}\\
$w_i$ & $\langle -6t^{3}+12t^{2}+48t+90,$\\
 & $\quad 24t^{2}+12t\rangle$\\
\addlinespace[2pt]
\multicolumn{2}{@{}p{0.97\textwidth}@{}}{$i=3$, $D_i=1$, $\rho_{ij}$: $7{:}2, \allowbreak 14{:}2, \allowbreak 15{:}1, \allowbreak 29{:}1, \allowbreak 30{:}1, \allowbreak 36{:}1, \allowbreak 39{:}2, \allowbreak 40{:}1, \allowbreak 41{:}2, \allowbreak 44{:}1$}\\
$w_i$ & $\langle -6t^{3}+9t^{2}-21t+18,$\\
 & $\quad -3t^{3}+3t^{2}-12t+12\rangle$\\
\addlinespace[2pt]
\multicolumn{2}{@{}p{0.97\textwidth}@{}}{$i=4$, $D_i=1$, $\rho_{ij}$: $7{:}2, \allowbreak 10{:}1, \allowbreak 11{:}1, \allowbreak 14{:}2, \allowbreak 27{:}2, \allowbreak 35{:}2, \allowbreak 41{:}1$}\\
$w_i$ & $\langle -2t^{3}+4t+12,$\\
 & $\quad 6t\rangle$\\
\addlinespace[2pt]
\multicolumn{2}{@{}p{0.97\textwidth}@{}}{$i=5$, $D_i=1$, $\rho_{ij}$: $1{:}2, \allowbreak 14{:}1, \allowbreak 32{:}2, \allowbreak 34{:}2, \allowbreak 35{:}2$}\\
$w_i$ & $\langle 3t^{3}-5t^{2}-12t+2,$\\
 & $\quad -7t^{3}-2t^{2}-23t+2\rangle$\\
\addlinespace[2pt]
\multicolumn{2}{@{}p{0.97\textwidth}@{}}{$i=6$, $D_i=1$, $\rho_{ij}$: $8{:}1, \allowbreak 18{:}2, \allowbreak 26{:}2, \allowbreak 37{:}2, \allowbreak 42{:}2, \allowbreak 43{:}1$}\\
$w_i$ & $\langle -44t^{3}+44t^{2}-20t-12,$\\
 & $\quad 4t^{3}+80t^{2}-44t+24\rangle$\\
\addlinespace[2pt]
\multicolumn{2}{@{}p{0.97\textwidth}@{}}{$i=7$, $D_i=1$, $\rho_{ij}$: $0{:}1, \allowbreak 1{:}2, \allowbreak 3{:}1, \allowbreak 4{:}1, \allowbreak 9{:}2, \allowbreak 17{:}1, \allowbreak 18{:}1, \allowbreak 19{:}1, \allowbreak 21{:}2, \allowbreak 24{:}1, \allowbreak 25{:}1, \allowbreak 31{:}1, \allowbreak 37{:}1, \allowbreak 38{:}1$}\\
$w_i$ & $\langle -11t^{5}+2t^{4}-14t^{3}+49t^{2}+14t+64,$\\
 & $\quad -4t^{5}+t^{4}-7t^{3}+23t^{2}+4t+32\rangle$\\
\addlinespace[2pt]
\multicolumn{2}{@{}p{0.97\textwidth}@{}}{$i=8$, $D_i=1$, $\rho_{ij}$: $0{:}1, \allowbreak 2{:}1, \allowbreak 6{:}2, \allowbreak 9{:}2, \allowbreak 13{:}1, \allowbreak 18{:}1, \allowbreak 24{:}2, \allowbreak 26{:}2, \allowbreak 34{:}1, \allowbreak 39{:}1, \allowbreak 40{:}1, \allowbreak 46{:}2$}\\
$w_i$ & $\langle 43t^{5}+105t^{4}+60t^{3}+175t^{2}+123t-60,$\\
 & $\quad 28t^{5}+133t^{4}+54t^{3}-74t^{2}-41t-162\rangle$\\
\addlinespace[2pt]
\multicolumn{2}{@{}p{0.97\textwidth}@{}}{$i=9$, $D_i=1$, $\rho_{ij}$: $7{:}1, \allowbreak 8{:}1, \allowbreak 25{:}1, \allowbreak 41{:}1$}\\
$w_i$ & $\langle 2t+2,$\\
 & $\quad 0\rangle$\\
\addlinespace[2pt]
\multicolumn{2}{@{}p{0.97\textwidth}@{}}{$i=10$, $D_i=1$, $\rho_{ij}$: $4{:}2, \allowbreak 19{:}2, \allowbreak 28{:}1, \allowbreak 29{:}2, \allowbreak 35{:}1, \allowbreak 39{:}2, \allowbreak 40{:}2, \allowbreak 42{:}1, \allowbreak 44{:}2$}\\
$w_i$ & $\langle -12t^{2}-26t-2,$\\
 & $\quad 6t^{2}-34t-4\rangle$\\
\addlinespace[2pt]
\multicolumn{2}{@{}p{0.97\textwidth}@{}}{$i=11$, $D_i=1$, $\rho_{ij}$: $0{:}2, \allowbreak 1{:}2, \allowbreak 4{:}2, \allowbreak 18{:}1, \allowbreak 24{:}2, \allowbreak 28{:}1, \allowbreak 35{:}1, \allowbreak 43{:}2$}\\
$w_i$ & $\langle 10t^{2}+38t+10,$\\
 & $\quad 26t^{2}+10t+2\rangle$\\
\addlinespace[2pt]
\multicolumn{2}{@{}p{0.97\textwidth}@{}}{$i=12$, $D_i=1$, $\rho_{ij}$: $22{:}2, \allowbreak 25{:}2, \allowbreak 29{:}2, \allowbreak 38{:}1, \allowbreak 44{:}1$}\\
$w_i$ & $\langle 2t^{2}-2t+4,$\\
 & $\quad 2t^{2}+2\rangle$\\
\addlinespace[2pt]
\multicolumn{2}{@{}p{0.97\textwidth}@{}}{$i=13$, $D_i=1$, $\rho_{ij}$: $0{:}2, \allowbreak 8{:}2, \allowbreak 17{:}1, \allowbreak 19{:}1, \allowbreak 24{:}1, \allowbreak 26{:}2, \allowbreak 32{:}2, \allowbreak 34{:}2, \allowbreak 35{:}2, \allowbreak 36{:}2, \allowbreak 39{:}1, \allowbreak 41{:}1, \allowbreak 46{:}2$}\\
$w_i$ & $\langle 36t^{2}-792t-348,$\\
 & $\quad 384t^{2}-324t-228\rangle$\\
\addlinespace[2pt]
\multicolumn{2}{@{}p{0.97\textwidth}@{}}{$i=14$, $D_i=1$, $\rho_{ij}$: $3{:}1, \allowbreak 4{:}1, \allowbreak 5{:}2, \allowbreak 17{:}1, \allowbreak 20{:}1, \allowbreak 22{:}1, \allowbreak 27{:}2, \allowbreak 29{:}2, \allowbreak 33{:}2, \allowbreak 46{:}2$}\\
$w_i$ & $\langle 12t^{2}-126t-6,$\\
 & $\quad 42t^{2}-42t\rangle$\\
\addlinespace[2pt]
\multicolumn{2}{@{}p{0.97\textwidth}@{}}{$i=15$, $D_i=1$, $\rho_{ij}$: $0{:}1, \allowbreak 2{:}2, \allowbreak 3{:}2, \allowbreak 22{:}2, \allowbreak 37{:}1, \allowbreak 39{:}2, \allowbreak 41{:}1$}\\
$w_i$ & $\langle 8t^{2}-24t-8,$\\
 & $\quad 4t^{2}-36t+44\rangle$\\
\addlinespace[2pt]
\multicolumn{2}{@{}p{0.97\textwidth}@{}}{$i=16$, $D_i=1$, $\rho_{ij}$: $0{:}2, \allowbreak 1{:}2, \allowbreak 23{:}2, \allowbreak 32{:}2, \allowbreak 33{:}2, \allowbreak 35{:}2, \allowbreak 38{:}1, \allowbreak 42{:}2$}\\
$w_i$ & $\langle 16t^{2}-108t+112,$\\
 & $\quad -14t^{2}+10t-98\rangle$\\
\addlinespace[2pt]
\multicolumn{2}{@{}p{0.97\textwidth}@{}}{$i=17$, $D_i=1$, $\rho_{ij}$: $7{:}2, \allowbreak 13{:}2, \allowbreak 14{:}2, \allowbreak 21{:}2, \allowbreak 22{:}1, \allowbreak 29{:}2, \allowbreak 31{:}1, \allowbreak 32{:}2, \allowbreak 36{:}2$}\\
$w_i$ & $\langle -6t^{3}+24t^{2}-42t,$\\
 & $\quad -18t^{3}+30t^{2}+6t-6\rangle$\\
\addlinespace[2pt]
\multicolumn{2}{@{}p{0.97\textwidth}@{}}{$i=18$, $D_i=1$, $\rho_{ij}$: $6{:}1, \allowbreak 7{:}2, \allowbreak 8{:}2, \allowbreak 11{:}2, \allowbreak 31{:}1, \allowbreak 37{:}1$}\\
$w_i$ & $\langle 8t^{3}+18t^{2}+18t+14,$\\
 & $\quad -2t^{3}+6t^{2}+6t+10\rangle$\\
\addlinespace[2pt]
\multicolumn{2}{@{}p{0.97\textwidth}@{}}{$i=19$, $D_i=1$, $\rho_{ij}$: $7{:}2, \allowbreak 10{:}1, \allowbreak 13{:}2, \allowbreak 32{:}1, \allowbreak 42{:}2$}\\
$w_i$ & $\langle 4t^{2}+6t+2,$\\
 & $\quad 2t^{2}+6t+4\rangle$\\
\addlinespace[2pt]
\multicolumn{2}{@{}p{0.97\textwidth}@{}}{$i=20$, $D_i=1$, $\rho_{ij}$: $0{:}1, \allowbreak 2{:}1, \allowbreak 14{:}2, \allowbreak 25{:}2, \allowbreak 26{:}1, \allowbreak 29{:}2, \allowbreak 30{:}2, \allowbreak 32{:}2, \allowbreak 36{:}2$}\\
$w_i$ & $\langle 0,$\\
 & $\quad -16t^{3}+12t^{2}+48t+16\rangle$\\
\addlinespace[2pt]
\multicolumn{2}{@{}p{0.97\textwidth}@{}}{$i=21$, $D_i=1$, $\rho_{ij}$: $7{:}1, \allowbreak 17{:}1, \allowbreak 25{:}1, \allowbreak 26{:}1, \allowbreak 28{:}2, \allowbreak 30{:}2, \allowbreak 35{:}2, \allowbreak 38{:}1, \allowbreak 43{:}2, \allowbreak 44{:}1, \allowbreak 45{:}1$}\\
$w_i$ & $\langle 132t^{3}+162t^{2}+198t,$\\
 & $\quad 6t^{3}-114t^{2}-144t-252\rangle$\\
\addlinespace[2pt]
\multicolumn{2}{@{}p{0.97\textwidth}@{}}{$i=22$, $D_i=1$, $\rho_{ij}$: $2{:}2, \allowbreak 12{:}1, \allowbreak 14{:}2, \allowbreak 15{:}1, \allowbreak 17{:}2, \allowbreak 41{:}2$}\\
$w_i$ & $\langle 4t^{3}-4t^{2}+4t,$\\
 & $\quad 4t^{3}-6t^{2}+6t-2\rangle$\\
\addlinespace[2pt]
\multicolumn{2}{@{}p{0.97\textwidth}@{}}{$i=23$, $D_i=1$, $\rho_{ij}$: $0{:}2, \allowbreak 16{:}1, \allowbreak 31{:}1, \allowbreak 32{:}2, \allowbreak 37{:}2, \allowbreak 40{:}2$}\\
$w_i$ & $\langle -4t^{3}+8t^{2}-8t+4,$\\
 & $\quad -8t^{3}+10t^{2}-10t+2\rangle$\\
\addlinespace[2pt]
\multicolumn{2}{@{}p{0.97\textwidth}@{}}{$i=24$, $D_i=1$, $\rho_{ij}$: $1{:}2, \allowbreak 2{:}1, \allowbreak 7{:}2, \allowbreak 8{:}1, \allowbreak 11{:}1, \allowbreak 13{:}2, \allowbreak 25{:}2, \allowbreak 35{:}2, \allowbreak 36{:}1, \allowbreak 39{:}2, \allowbreak 44{:}1$}\\
$w_i$ & $\langle -252t^{2}+576t-432,$\\
 & $\quad 36t^{3}-156t^{2}+216t-96\rangle$\\
\addlinespace[2pt]
\multicolumn{2}{@{}p{0.97\textwidth}@{}}{$i=25$, $D_i=1$, $\rho_{ij}$: $0{:}2, \allowbreak 7{:}2, \allowbreak 9{:}2, \allowbreak 12{:}1, \allowbreak 20{:}1, \allowbreak 21{:}2, \allowbreak 24{:}1, \allowbreak 38{:}1, \allowbreak 44{:}2$}\\
$w_i$ & $\langle -14t^{3}-48t^{2}+2t+24,$\\
 & $\quad -24t^{3}+42t+24\rangle$\\
\addlinespace[2pt]
\multicolumn{2}{@{}p{0.97\textwidth}@{}}{$i=26$, $D_i=1$, $\rho_{ij}$: $6{:}1, \allowbreak 8{:}1, \allowbreak 13{:}1, \allowbreak 20{:}2, \allowbreak 21{:}2, \allowbreak 29{:}1$}\\
$w_i$ & $\langle -t^{3}-4t^{2}+3t+14,$\\
 & $\quad -5t^{3}-13t^{2}-6t\rangle$\\
\addlinespace[2pt]
\multicolumn{2}{@{}p{0.97\textwidth}@{}}{$i=27$, $D_i=1$, $\rho_{ij}$: $1{:}1, \allowbreak 2{:}2, \allowbreak 4{:}1, \allowbreak 14{:}1, \allowbreak 45{:}2$}\\
$w_i$ & $\langle -t^{2}+t+2,$\\
 & $\quad t^{3}+t^{2}+2t+2\rangle$\\
\addlinespace[2pt]
\multicolumn{2}{@{}p{0.97\textwidth}@{}}{$i=28$, $D_i=1$, $\rho_{ij}$: $1{:}2, \allowbreak 10{:}2, \allowbreak 11{:}2, \allowbreak 21{:}1, \allowbreak 43{:}2$}\\
$w_i$ & $\langle 2t^{3}+4t^{2}-2t-4,$\\
 & $\quad -2t^{3}-10t^{2}-16t-8\rangle$\\
\addlinespace[2pt]
\multicolumn{2}{@{}p{0.97\textwidth}@{}}{$i=29$, $D_i=1$, $\rho_{ij}$: $0{:}2, \allowbreak 2{:}2, \allowbreak 3{:}2, \allowbreak 10{:}1, \allowbreak 12{:}1, \allowbreak 14{:}1, \allowbreak 17{:}1, \allowbreak 20{:}1, \allowbreak 26{:}2, \allowbreak 36{:}1, \allowbreak 38{:}2, \allowbreak 44{:}2$}\\
$w_i$ & $\langle 354t^{3}-658t^{2}+220t+248,$\\
 & $\quad 180t^{3}-146t^{2}+50t-404\rangle$\\
\addlinespace[2pt]
\multicolumn{2}{@{}p{0.97\textwidth}@{}}{$i=30$, $D_i=1$, $\rho_{ij}$: $2{:}2, \allowbreak 3{:}2, \allowbreak 20{:}1, \allowbreak 21{:}1, \allowbreak 32{:}1, \allowbreak 36{:}1, \allowbreak 37{:}2$}\\
$w_i$ & $\langle -6t^{3}+12t^{2}-6t,$\\
 & $\quad -3t^{3}+12t^{2}-15t+6\rangle$\\
\addlinespace[2pt]
\multicolumn{2}{@{}p{0.97\textwidth}@{}}{$i=31$, $D_i=1$, $\rho_{ij}$: $7{:}2, \allowbreak 17{:}2, \allowbreak 18{:}2, \allowbreak 23{:}2, \allowbreak 40{:}2$}\\
$w_i$ & $\langle 4t^{2},$\\
 & $\quad 0\rangle$\\
\addlinespace[2pt]
\multicolumn{2}{@{}p{0.97\textwidth}@{}}{$i=32$, $D_i=1$, $\rho_{ij}$: $5{:}1, \allowbreak 13{:}1, \allowbreak 16{:}1, \allowbreak 17{:}1, \allowbreak 19{:}2, \allowbreak 20{:}1, \allowbreak 23{:}1, \allowbreak 30{:}2, \allowbreak 35{:}1, \allowbreak 41{:}1, \allowbreak 43{:}1, \allowbreak 45{:}1, \allowbreak 46{:}1$}\\
$w_i$ & $\langle -44t^{3}-128t^{2}-160t+32,$\\
 & $\quad 20t^{3}-4t^{2}-176t-80\rangle$\\
\addlinespace[2pt]
\multicolumn{2}{@{}p{0.97\textwidth}@{}}{$i=33$, $D_i=1$, $\rho_{ij}$: $1{:}2, \allowbreak 2{:}1, \allowbreak 14{:}1, \allowbreak 16{:}1, \allowbreak 35{:}2, \allowbreak 38{:}2, \allowbreak 40{:}1, \allowbreak 43{:}2$}\\
$w_i$ & $\langle 4t^{2}+8t+4,$\\
 & $\quad 6t^{3}-4t+2\rangle$\\
\addlinespace[2pt]
\multicolumn{2}{@{}p{0.97\textwidth}@{}}{$i=34$, $D_i=1$, $\rho_{ij}$: $5{:}1, \allowbreak 8{:}2, \allowbreak 13{:}1, \allowbreak 41{:}2, \allowbreak 43{:}2$}\\
$w_i$ & $\langle 0,$\\
 & $\quad 2t^{3}+2t^{2}+6t-2\rangle$\\
\addlinespace[2pt]
\multicolumn{2}{@{}p{0.97\textwidth}@{}}{$i=35$, $D_i=1$, $\rho_{ij}$: $4{:}1, \allowbreak 5{:}1, \allowbreak 10{:}2, \allowbreak 11{:}2, \allowbreak 13{:}1, \allowbreak 16{:}1, \allowbreak 21{:}1, \allowbreak 24{:}1, \allowbreak 32{:}2, \allowbreak 33{:}1, \allowbreak 39{:}2, \allowbreak 41{:}2, \allowbreak 42{:}2$}\\
$w_i$ & $\langle -32t^{4}-8t^{3}-58t^{2}+68t+54,$\\
 & $\quad -6t^{5}-28t^{4}+32t^{3}+4t^{2}+58t+60\rangle$\\
\addlinespace[2pt]
\multicolumn{2}{@{}p{0.97\textwidth}@{}}{$i=36$, $D_i=1$, $\rho_{ij}$: $3{:}2, \allowbreak 13{:}1, \allowbreak 17{:}1, \allowbreak 20{:}1, \allowbreak 24{:}2, \allowbreak 29{:}2, \allowbreak 30{:}2, \allowbreak 37{:}1$}\\
$w_i$ & $\langle -3t^{5}-12t^{4}-3t^{3}-48t^{2}-30t-30,$\\
 & $\quad 0\rangle$\\
\addlinespace[2pt]
\multicolumn{2}{@{}p{0.97\textwidth}@{}}{$i=37$, $D_i=1$, $\rho_{ij}$: $6{:}1, \allowbreak 7{:}2, \allowbreak 15{:}2, \allowbreak 18{:}2, \allowbreak 23{:}1, \allowbreak 30{:}1, \allowbreak 36{:}2$}\\
$w_i$ & $\langle 6t^{2}-9t+3,$\\
 & $\quad 3t^{3}-3t^{2}-3t+3\rangle$\\
\addlinespace[2pt]
\multicolumn{2}{@{}p{0.97\textwidth}@{}}{$i=38$, $D_i=1$, $\rho_{ij}$: $7{:}2, \allowbreak 12{:}2, \allowbreak 16{:}2, \allowbreak 21{:}2, \allowbreak 25{:}2, \allowbreak 29{:}1, \allowbreak 33{:}1, \allowbreak 40{:}1$}\\
$w_i$ & $\langle 2t^{3}+2t+2,$\\
 & $\quad 4t^{3}-2t^{2}+6t-8\rangle$\\
\addlinespace[2pt]
\multicolumn{2}{@{}p{0.97\textwidth}@{}}{$i=39$, $D_i=1$, $\rho_{ij}$: $3{:}1, \allowbreak 8{:}2, \allowbreak 10{:}1, \allowbreak 13{:}2, \allowbreak 15{:}1, \allowbreak 24{:}1, \allowbreak 35{:}1, \allowbreak 46{:}1$}\\
$w_i$ & $\langle -18t^{2}-42t,$\\
 & $\quad -12t^{3}-24t^{2}+24\rangle$\\
\addlinespace[2pt]
\multicolumn{2}{@{}p{0.97\textwidth}@{}}{$i=40$, $D_i=1$, $\rho_{ij}$: $0{:}1, \allowbreak 2{:}2, \allowbreak 3{:}2, \allowbreak 8{:}2, \allowbreak 10{:}1, \allowbreak 23{:}1, \allowbreak 31{:}1, \allowbreak 33{:}2, \allowbreak 38{:}2, \allowbreak 45{:}2$}\\
$w_i$ & $\langle 23t^{3}+37t^{2}-84t+30,$\\
 & $\quad -12t^{3}+75t^{2}-57t-48\rangle$\\
\addlinespace[2pt]
\multicolumn{2}{@{}p{0.97\textwidth}@{}}{$i=41$, $D_i=1$, $\rho_{ij}$: $2{:}1, \allowbreak 3{:}1, \allowbreak 4{:}2, \allowbreak 9{:}2, \allowbreak 13{:}2, \allowbreak 15{:}2, \allowbreak 22{:}1, \allowbreak 32{:}2, \allowbreak 34{:}1, \allowbreak 35{:}1$}\\
$w_i$ & $\langle -2t^{3}+43t^{2}-35t-12,$\\
 & $\quad 13t^{3}+15t^{2}-40t+30\rangle$\\
\addlinespace[2pt]
\multicolumn{2}{@{}p{0.97\textwidth}@{}}{$i=42$, $D_i=1$, $\rho_{ij}$: $2{:}1, \allowbreak 6{:}1, \allowbreak 10{:}2, \allowbreak 16{:}1, \allowbreak 19{:}1, \allowbreak 35{:}1, \allowbreak 45{:}1$}\\
$w_i$ & $\langle 2t^{3}+2t^{2}-6t-6,$\\
 & $\quad -6t-6\rangle$\\
\addlinespace[2pt]
\multicolumn{2}{@{}p{0.97\textwidth}@{}}{$i=43$, $D_i=1$, $\rho_{ij}$: $0{:}2, \allowbreak 1{:}2, \allowbreak 2{:}2, \allowbreak 6{:}2, \allowbreak 11{:}1, \allowbreak 21{:}1, \allowbreak 28{:}1, \allowbreak 32{:}2, \allowbreak 33{:}1, \allowbreak 34{:}1, \allowbreak 45{:}1$}\\
$w_i$ & $\langle -6t^{5}-28t^{4}-6t^{3}+24t^{2}+14t+30,$\\
 & $\quad -11t^{5}-24t^{4}-6t^{3}-23t^{2}-12t+18\rangle$\\
\addlinespace[2pt]
\multicolumn{2}{@{}p{0.97\textwidth}@{}}{$i=44$, $D_i=1$, $\rho_{ij}$: $3{:}2, \allowbreak 10{:}1, \allowbreak 12{:}2, \allowbreak 21{:}2, \allowbreak 24{:}2, \allowbreak 25{:}1, \allowbreak 29{:}1$}\\
$w_i$ & $\langle -4t^{3}+4t^{2}+6t-2,$\\
 & $\quad -2t^{3}+6t^{2}+2t-6\rangle$\\
\addlinespace[2pt]
\multicolumn{2}{@{}p{0.97\textwidth}@{}}{$i=45$, $D_i=1$, $\rho_{ij}$: $0{:}2, \allowbreak 21{:}2, \allowbreak 27{:}1, \allowbreak 32{:}2, \allowbreak 40{:}1, \allowbreak 42{:}2, \allowbreak 43{:}2$}\\
$w_i$ & $\langle -48t+84,$\\
 & $\quad -36t+24\rangle$\\
\addlinespace[2pt]
\multicolumn{2}{@{}p{0.97\textwidth}@{}}{$i=46$, $D_i=1$, $\rho_{ij}$: $8{:}1, \allowbreak 13{:}1, \allowbreak 14{:}1, \allowbreak 32{:}2, \allowbreak 39{:}2$}\\
$w_i$ & $\langle 6t^{2}+18t+18,$\\
 & $\quad 6t^{2}+18t+18\rangle$\\
\addlinespace[2pt]
\end{longtable}
\endgroup

\begingroup\small
\setlength{\LTcapwidth}{\textwidth}
\begin{longtable}{@{}lrl@{\qquad}lrl@{}}
\caption{The edges $(i,j)$ with a linear endpoint, the linear endpoint $k$ used, and the norm of $R_{ij}=\operatorname{Res}_t(f_k,g_{kl})$.}\label{tab:sep}\\
\hline $(i,j)$ & $k$ & $N(R_{ij})$ & $(i,j)$ & $k$ & $N(R_{ij})$\\\hline\endfirsthead\hline $(i,j)$ & $k$ & $N(R_{ij})$ & $(i,j)$ & $k$ & $N(R_{ij})$\\\hline\endhead\hline\endfoot
$(0,7)$ & 0 & $2^{6}\cdot 3^{2}$ & $(12,29)$ & 12 & $2^{2}\cdot 13$\\
$(0,8)$ & 0 & $3^{3}\cdot 13$ & $(12,38)$ & 12 & $2^{2}\cdot 3$\\
$(0,11)$ & 0 & $2^{2}\cdot 3^{2}$ & $(12,44)$ & 12 & $2^{2}\cdot 13$\\
$(0,13)$ & 0 & $2^{2}\cdot 13$ & $(13,17)$ & 13 & $2^{2}\cdot 7$\\
$(0,15)$ & 0 & $2^{2}\cdot 3$ & $(13,19)$ & 13 & $2^{4}\cdot 3^{2}$\\
$(0,16)$ & 0 & $2^{6}\cdot 3$ & $(13,24)$ & 13 & $3^{2}$\\
$(0,20)$ & 0 & $2^{2}\cdot 13$ & $(13,26)$ & 13 & $2^{4}$\\
$(0,23)$ & 0 & $2^{2}\cdot 3^{2}\cdot 13$ & $(13,32)$ & 13 & $2^{2}\cdot 3^{2}\cdot 7$\\
$(0,25)$ & 0 & $2^{6}\cdot 3$ & $(13,34)$ & 13 & $2^{4}$\\
$(0,29)$ & 0 & $2^{4}\cdot 3^{2}\cdot 13$ & $(13,35)$ & 13 & $2^{4}\cdot 3^{2}$\\
$(0,40)$ & 0 & $2^{2}\cdot 3\cdot 13$ & $(13,36)$ & 13 & $2^{2}\cdot 3^{2}\cdot 7^{2}$\\
$(0,43)$ & 0 & $2^{4}\cdot 3^{3}$ & $(13,39)$ & 13 & $2^{2}\cdot 3^{2}\cdot 7$\\
$(0,45)$ & 0 & $2^{4}\cdot 3^{2}$ & $(13,41)$ & 13 & $2^{2}\cdot 3^{2}$\\
$(1,5)$ & 1 & $2^{8}\cdot 7$ & $(13,46)$ & 13 & $2^{2}\cdot 3^{2}\cdot 7^{2}$\\
$(1,7)$ & 1 & $2^{2}\cdot 7$ & $(14,17)$ & 14 & $2^{2}\cdot 3^{2}$\\
$(1,11)$ & 1 & $2^{2}\cdot 7$ & $(14,20)$ & 14 & $2^{4}\cdot 13$\\
$(1,16)$ & 1 & $2^{4}\cdot 13$ & $(14,22)$ & 14 & $2^{2}\cdot 3^{2}$\\
$(1,24)$ & 1 & $2^{2}\cdot 7$ & $(14,27)$ & 14 & $3^{2}\cdot 7^{2}$\\
$(1,27)$ & 1 & $2^{2}\cdot 7^{2}\cdot 13$ & $(14,29)$ & 14 & $3^{2}\cdot 13$\\
$(1,28)$ & 1 & $2^{2}\cdot 7^{2}$ & $(14,33)$ & 14 & $2^{6}$\\
$(1,33)$ & 1 & $2^{2}\cdot 7$ & $(14,46)$ & 14 & $2^{4}\cdot 3^{2}\cdot 7$\\
$(1,43)$ & 1 & $7$ & $(15,22)$ & 15 & $2^{4}\cdot 13$\\
$(2,8)$ & 2 & $2^{4}\cdot 3^{3}$ & $(15,37)$ & 15 & $2^{4}$\\
$(2,15)$ & 2 & $2^{6}\cdot 3$ & $(15,39)$ & 15 & $2^{2}\cdot 13^{2}$\\
$(2,20)$ & 2 & $2^{2}\cdot 13$ & $(15,41)$ & 15 & $2^{2}\cdot 3\cdot 13$\\
$(2,22)$ & 2 & $2^{2}\cdot 3^{2}\cdot 13$ & $(16,23)$ & 16 & $2^{4}\cdot 13$\\
$(2,24)$ & 2 & $2^{2}\cdot 3^{2}$ & $(16,32)$ & 16 & $2^{2}\cdot 13^{2}$\\
$(2,27)$ & 2 & $2^{2}\cdot 3^{2}$ & $(16,33)$ & 16 & $2^{2}\cdot 13$\\
$(2,29)$ & 2 & $2^{2}\cdot 3^{2}$ & $(16,35)$ & 16 & $2^{2}\cdot 3$\\
$(2,30)$ & 2 & $2^{4}\cdot 3^{2}\cdot 13$ & $(16,38)$ & 16 & $2^{4}$\\
$(2,33)$ & 2 & $2^{2}\cdot 13$ & $(16,42)$ & 16 & $2^{2}\cdot 3\cdot 13$\\
$(2,40)$ & 2 & $2^{2}\cdot 3^{2}$ & $(17,21)$ & 17 & $2^{2}\cdot 3^{2}$\\
$(2,41)$ & 2 & $2^{2}\cdot 3\cdot 13$ & $(17,22)$ & 17 & $2^{2}\cdot 3^{2}$\\
$(2,42)$ & 2 & $2^{2}\cdot 3$ & $(17,29)$ & 17 & $2^{2}\cdot 3^{4}\cdot 7$\\
$(2,43)$ & 2 & $3^{3}\cdot 13$ & $(17,31)$ & 17 & $2^{4}$\\
$(3,7)$ & 3 & $2^{2}\cdot 3^{2}$ & $(17,32)$ & 17 & $2^{2}\cdot 3^{2}$\\
$(3,14)$ & 3 & $2^{2}\cdot 3^{2}$ & $(17,36)$ & 17 & $3^{5}\cdot 7$\\
$(3,15)$ & 3 & $2^{2}\cdot 3\cdot 13$ & $(18,31)$ & 18 & $3$\\
$(3,29)$ & 3 & $2^{4}\cdot 3^{2}$ & $(18,37)$ & 18 & $2^{2}\cdot 3^{2}$\\
$(3,30)$ & 3 & $2^{4}\cdot 3^{2}$ & $(19,32)$ & 19 & $2^{8}\cdot 3$\\
$(3,36)$ & 3 & $2^{4}\cdot 3^{4}$ & $(19,42)$ & 19 & $2^{2}\cdot 3$\\
$(3,39)$ & 3 & $2^{2}\cdot 3^{2}\cdot 13$ & $(20,25)$ & 20 & $2^{4}\cdot 13$\\
$(3,40)$ & 3 & $2^{2}\cdot 3$ & $(20,26)$ & 20 & $2^{6}$\\
$(3,41)$ & 3 & $2^{2}\cdot 3$ & $(20,29)$ & 20 & $2^{6}\cdot 13^{2}$\\
$(3,44)$ & 3 & $2^{4}$ & $(20,30)$ & 20 & $2^{6}\cdot 13^{2}$\\
$(4,7)$ & 4 & $2^{2}\cdot 7$ & $(20,32)$ & 20 & $2^{6}$\\
$(4,10)$ & 4 & $2^{2}\cdot 7$ & $(20,36)$ & 20 & $2^{4}\cdot 13^{2}$\\
$(4,11)$ & 4 & $2^{2}\cdot 7$ & $(21,25)$ & 21 & $2^{2}\cdot 13$\\
$(4,14)$ & 4 & $2^{2}\cdot 7$ & $(21,26)$ & 21 & $2^{6}\cdot 7$\\
$(4,27)$ & 4 & $2^{2}\cdot 7^{2}$ & $(21,28)$ & 21 & $2^{4}\cdot 3^{2}$\\
$(4,35)$ & 4 & $7\cdot 13$ & $(21,30)$ & 21 & $3^{2}\cdot 13$\\
$(4,41)$ & 4 & $2^{4}\cdot 13$ & $(21,35)$ & 21 & $2^{4}\cdot 3^{3}$\\
$(5,14)$ & 5 & $2^{2}\cdot 7$ & $(21,38)$ & 21 & $2^{2}\cdot 3^{4}\cdot 7$\\
$(5,32)$ & 5 & $3^{3}\cdot 7^{3}$ & $(21,43)$ & 21 & $3^{3}\cdot 13$\\
$(5,34)$ & 5 & $1$ & $(21,44)$ & 21 & $2^{2}\cdot 7$\\
$(5,35)$ & 5 & $2^{6}\cdot 3^{3}$ & $(21,45)$ & 21 & $2^{2}\cdot 3^{2}\cdot 13$\\
$(6,8)$ & 6 & $2^{2}\cdot 3^{3}\cdot 13$ & $(22,41)$ & 22 & $3\cdot 13^{3}$\\
$(6,18)$ & 6 & $2^{2}\cdot 3^{2}\cdot 13$ & $(23,31)$ & 23 & $2^{4}\cdot 3$\\
$(6,26)$ & 6 & $2^{4}\cdot 7^{4}$ & $(23,32)$ & 23 & $2^{2}\cdot 3^{6}$\\
$(6,37)$ & 6 & $2^{2}\cdot 3^{2}$ & $(23,37)$ & 23 & $2^{2}\cdot 3^{2}$\\
$(6,42)$ & 6 & $2^{10}\cdot 3^{2}$ & $(23,40)$ & 23 & $3\cdot 13^{3}$\\
$(6,43)$ & 6 & $3^{3}$ & $(24,25)$ & 24 & $2^{8}\cdot 7$\\
$(7,9)$ & 7 & $2^{6}$ & $(24,35)$ & 24 & $2^{4}\cdot 3^{3}$\\
$(7,17)$ & 7 & $3^{5}$ & $(24,36)$ & 24 & $2^{4}\cdot 3^{6}$\\
$(7,18)$ & 7 & $3^{2}$ & $(24,39)$ & 24 & $2^{6}\cdot 3^{4}$\\
$(7,19)$ & 7 & $2^{4}\cdot 3^{2}\cdot 7$ & $(24,44)$ & 24 & $2^{2}\cdot 7$\\
$(7,21)$ & 7 & $3^{2}$ & $(25,38)$ & 25 & $2^{2}\cdot 7$\\
$(7,24)$ & 7 & $2^{6}\cdot 3^{2}\cdot 7^{2}$ & $(25,44)$ & 25 & $2^{8}$\\
$(7,25)$ & 7 & $7^{2}$ & $(26,29)$ & 26 & $2^{10}\cdot 7$\\
$(7,31)$ & 7 & $2^{4}$ & $(27,45)$ & 27 & $2^{2}\cdot 3^{2}\cdot 13$\\
$(7,37)$ & 7 & $3^{5}$ & $(28,43)$ & 28 & $2^{6}\cdot 3^{2}$\\
$(7,38)$ & 7 & $2^{2}\cdot 7$ & $(29,36)$ & 29 & $2^{2}\cdot 3^{4}\cdot 7^{2}\cdot 13$\\
$(8,9)$ & 8 & $2^{2}\cdot 3^{2}$ & $(29,38)$ & 29 & $2^{2}\cdot 7$\\
$(8,13)$ & 8 & $2^{2}\cdot 13^{2}$ & $(29,44)$ & 29 & $7\cdot 13$\\
$(8,18)$ & 8 & $3^{3}\cdot 13$ & $(30,32)$ & 30 & $2^{2}\cdot 3^{2}\cdot 7^{2}$\\
$(8,24)$ & 8 & $2^{6}\cdot 3^{4}$ & $(30,36)$ & 30 & $2^{2}\cdot 3^{4}\cdot 7^{2}\cdot 13$\\
$(8,26)$ & 8 & $2^{12}$ & $(30,37)$ & 30 & $2^{2}\cdot 3^{4}\cdot 7$\\
$(8,34)$ & 8 & $2^{6}$ & $(31,40)$ & 31 & $3$\\
$(8,39)$ & 8 & $2^{2}\cdot 3^{3}\cdot 7^{2}$ & $(32,35)$ & 32 & $2^{4}\cdot 3^{5}$\\
$(8,40)$ & 8 & $2^{2}\cdot 3\cdot 7^{2}$ & $(32,41)$ & 32 & $3^{3}\cdot 7^{3}$\\
$(8,46)$ & 8 & $2^{6}\cdot 3^{4}\cdot 7$ & $(32,43)$ & 32 & $2^{2}\cdot 3^{3}\cdot 7^{2}$\\
$(9,25)$ & 9 & $2^{2}\cdot 3^{2}$ & $(32,45)$ & 32 & $2^{2}\cdot 3^{2}\cdot 13$\\
$(9,41)$ & 9 & $2^{2}\cdot 3^{2}$ & $(32,46)$ & 32 & $2^{2}\cdot 3^{4}\cdot 7^{2}$\\
$(10,19)$ & 10 & $2^{2}\cdot 7$ & $(33,35)$ & 33 & $2^{4}\cdot 7$\\
$(10,28)$ & 10 & $2^{2}\cdot 3^{2}$ & $(33,38)$ & 33 & $2^{2}\cdot 7^{2}$\\
$(10,29)$ & 10 & $2^{2}\cdot 3^{2}$ & $(33,40)$ & 33 & $2^{2}\cdot 7$\\
$(10,35)$ & 10 & $2^{2}\cdot 3^{2}$ & $(33,43)$ & 33 & $2^{2}\cdot 13^{2}$\\
$(10,39)$ & 10 & $2^{2}\cdot 3^{2}$ & $(34,41)$ & 34 & $2^{6}$\\
$(10,40)$ & 10 & $2^{6}$ & $(34,43)$ & 34 & $2^{6}$\\
$(10,42)$ & 10 & $1$ & $(35,39)$ & 35 & $2^{4}\cdot 3^{3}\cdot 7$\\
$(10,44)$ & 10 & $2^{2}\cdot 7$ & $(35,41)$ & 35 & $3^{5}\cdot 13$\\
$(11,18)$ & 11 & $2^{2}\cdot 3^{2}$ & $(35,42)$ & 35 & $2^{10}\cdot 3$\\
$(11,24)$ & 11 & $2^{2}\cdot 3^{2}$ & $(36,37)$ & 36 & $3^{5}\cdot 7$\\
$(11,28)$ & 11 & $2^{4}\cdot 3^{2}\cdot 7$ & $(38,40)$ & 38 & $2^{2}\cdot 3^{2}$\\
$(11,35)$ & 11 & $2^{4}\cdot 3^{2}\cdot 7$ & $(39,46)$ & 39 & $2^{2}\cdot 3^{4}\cdot 7^{2}$\\
$(11,43)$ & 11 & $2^{2}\cdot 3^{2}$ & $(40,45)$ & 40 & $2^{4}\cdot 13$\\
$(12,22)$ & 12 & $2^{4}\cdot 3$ & $(42,45)$ & 42 & $2^{2}\cdot 13^{2}$\\
$(12,25)$ & 12 & $2^{4}$ & $(43,45)$ & 43 & $2^{2}\cdot 3^{2}\cdot 13$\\
\end{longtable}
\endgroup

\begingroup\small
\setlength{\LTcapwidth}{\textwidth}
\begin{longtable}{@{}rrl>{\raggedright\arraybackslash}p{0.66\textwidth}@{}}
\caption{The terminal residue classes $x\equiv a$, $y\equiv b\pmod{p^k}$ of Step~4, grouped by $p$, $k$ and the value of $B_p$.}\label{tab:local}\\
\hline $p$ & $k$ & $B_p$ & classes $(a,b)$\\\hline\endfirsthead
\hline $p$ & $k$ & $B_p$ & classes $(a,b)$\\\hline\endhead\hline\endfoot
2 & 3 & $0$ & $(1,3)$,\allowbreak\ $(2,3)$,\allowbreak\ $(3,5)$,\allowbreak\ $(4,5)$,\allowbreak\ $(5,7)$\\
2 & 4 & $0$ & $(0,9)$,\allowbreak\ $(6,15)$,\allowbreak\ $(7,1)$,\allowbreak\ $(8,1)$,\allowbreak\ $(14,7)$,\allowbreak\ $(15,9)$\\
\addlinespace[3pt]
3 & 3 & $0$ & $(4,6)$,\allowbreak\ $(4,15)$,\allowbreak\ $(4,24)$,\allowbreak\ $(6,2)$,\allowbreak\ $(6,5)$,\allowbreak\ $(6,8)$,\allowbreak\ $(6,11)$,\allowbreak\ $(6,14)$,\allowbreak\ $(6,17)$,\allowbreak\ $(6,20)$,\allowbreak\ $(6,23)$,\allowbreak\ $(6,26)$,\allowbreak\ $(13,3)$,\allowbreak\ $(13,12)$,\allowbreak\ $(13,21)$,\allowbreak\ $(22,0)$,\allowbreak\ $(22,9)$,\allowbreak\ $(22,18)$\\
\addlinespace[3pt]
7 & 2 & $1$ & $(0,9)$,\allowbreak\ $(1,25)$,\allowbreak\ $(3,37)$,\allowbreak\ $(5,23)$,\allowbreak\ $(6,2)$,\allowbreak\ $(7,44)$,\allowbreak\ $(8,18)$,\allowbreak\ $(12,16)$,\allowbreak\ $(13,44)$,\allowbreak\ $(14,30)$,\allowbreak\ $(15,11)$,\allowbreak\ $(17,23)$,\allowbreak\ $(19,9)$,\allowbreak\ $(21,16)$,\allowbreak\ $(22,4)$,\allowbreak\ $(26,2)$,\allowbreak\ $(27,30)$,\allowbreak\ $(28,2)$,\allowbreak\ $(29,46)$,\allowbreak\ $(31,9)$,\allowbreak\ $(33,44)$,\allowbreak\ $(34,23)$,\allowbreak\ $(35,37)$,\allowbreak\ $(38,2)$,\allowbreak\ $(40,37)$,\allowbreak\ $(42,23)$,\allowbreak\ $(43,32)$,\allowbreak\ $(45,44)$,\allowbreak\ $(48,9)$\\
7 & 3 & $1$ & $(10,79)$,\allowbreak\ $(20,135)$,\allowbreak\ $(24,310)$,\allowbreak\ $(36,137)$,\allowbreak\ $(41,212)$,\allowbreak\ $(47,128)$,\allowbreak\ $(59,30)$,\allowbreak\ $(69,86)$,\allowbreak\ $(73,261)$,\allowbreak\ $(85,88)$,\allowbreak\ $(90,163)$,\allowbreak\ $(96,79)$,\allowbreak\ $(108,324)$,\allowbreak\ $(118,37)$,\allowbreak\ $(122,212)$,\allowbreak\ $(134,39)$,\allowbreak\ $(139,114)$,\allowbreak\ $(145,30)$,\allowbreak\ $(157,275)$,\allowbreak\ $(167,331)$,\allowbreak\ $(183,333)$,\allowbreak\ $(188,65)$,\allowbreak\ $(194,324)$,\allowbreak\ $(206,226)$,\allowbreak\ $(216,282)$,\allowbreak\ $(220,114)$,\allowbreak\ $(232,284)$,\allowbreak\ $(237,16)$,\allowbreak\ $(243,275)$,\allowbreak\ $(255,177)$,\allowbreak\ $(265,233)$,\allowbreak\ $(269,65)$,\allowbreak\ $(281,235)$,\allowbreak\ $(286,310)$,\allowbreak\ $(292,226)$,\allowbreak\ $(304,128)$,\allowbreak\ $(314,184)$,\allowbreak\ $(318,16)$,\allowbreak\ $(330,186)$,\allowbreak\ $(335,261)$,\allowbreak\ $(341,177)$\\
7 & 4 & $1$ & $(171,849)$,\allowbreak\ $(514,506)$,\allowbreak\ $(857,163)$,\allowbreak\ $(1200,2221)$,\allowbreak\ $(1543,1878)$,\allowbreak\ $(1886,1535)$,\allowbreak\ $(2229,1192)$\\
\addlinespace[3pt]
13 & 1 & $0$ & $(0,5)$,\allowbreak\ $(2,1)$,\allowbreak\ $(2,2)$,\allowbreak\ $(3,3)$,\allowbreak\ $(3,11)$,\allowbreak\ $(3,12)$,\allowbreak\ $(6,4)$,\allowbreak\ $(8,9)$,\allowbreak\ $(9,1)$,\allowbreak\ $(9,2)$,\allowbreak\ $(10,5)$,\allowbreak\ $(12,5)$\\
13 & 2 & $0$ & $(2,49)$,\allowbreak\ $(4,148)$,\allowbreak\ $(5,3)$,\allowbreak\ $(5,11)$,\allowbreak\ $(5,155)$,\allowbreak\ $(9,127)$,\allowbreak\ $(14,95)$,\allowbreak\ $(15,36)$,\allowbreak\ $(17,161)$,\allowbreak\ $(18,51)$,\allowbreak\ $(18,128)$,\allowbreak\ $(18,159)$,\allowbreak\ $(22,166)$,\allowbreak\ $(27,147)$,\allowbreak\ $(28,23)$,\allowbreak\ $(30,5)$,\allowbreak\ $(31,76)$,\allowbreak\ $(31,116)$,\allowbreak\ $(31,146)$,\allowbreak\ $(35,36)$,\allowbreak\ $(40,30)$,\allowbreak\ $(41,10)$,\allowbreak\ $(43,18)$,\allowbreak\ $(44,12)$,\allowbreak\ $(44,133)$,\allowbreak\ $(48,75)$,\allowbreak\ $(53,82)$,\allowbreak\ $(54,166)$,\allowbreak\ $(56,31)$,\allowbreak\ $(57,77)$,\allowbreak\ $(57,120)$,\allowbreak\ $(57,141)$,\allowbreak\ $(61,114)$,\allowbreak\ $(66,134)$,\allowbreak\ $(67,153)$,\allowbreak\ $(69,44)$,\allowbreak\ $(70,89)$,\allowbreak\ $(70,107)$,\allowbreak\ $(70,142)$,\allowbreak\ $(74,153)$,\allowbreak\ $(79,17)$,\allowbreak\ $(80,140)$,\allowbreak\ $(82,57)$,\allowbreak\ $(83,37)$,\allowbreak\ $(83,38)$,\allowbreak\ $(83,94)$,\allowbreak\ $(87,23)$,\allowbreak\ $(92,69)$,\allowbreak\ $(93,127)$,\allowbreak\ $(95,70)$,\allowbreak\ $(96,81)$,\allowbreak\ $(96,103)$,\allowbreak\ $(96,154)$,\allowbreak\ $(100,62)$,\allowbreak\ $(105,121)$,\allowbreak\ $(106,114)$,\allowbreak\ $(108,83)$,\allowbreak\ $(109,68)$,\allowbreak\ $(109,102)$,\allowbreak\ $(109,168)$,\allowbreak\ $(113,101)$,\allowbreak\ $(118,4)$,\allowbreak\ $(119,101)$,\allowbreak\ $(121,96)$,\allowbreak\ $(122,50)$,\allowbreak\ $(122,55)$,\allowbreak\ $(122,64)$,\allowbreak\ $(126,140)$,\allowbreak\ $(131,56)$,\allowbreak\ $(132,88)$,\allowbreak\ $(134,109)$,\allowbreak\ $(135,42)$,\allowbreak\ $(135,129)$,\allowbreak\ $(135,167)$,\allowbreak\ $(139,10)$,\allowbreak\ $(144,108)$,\allowbreak\ $(145,75)$,\allowbreak\ $(147,122)$,\allowbreak\ $(148,29)$,\allowbreak\ $(148,115)$,\allowbreak\ $(152,49)$,\allowbreak\ $(157,160)$,\allowbreak\ $(160,135)$,\allowbreak\ $(161,16)$,\allowbreak\ $(161,63)$,\allowbreak\ $(161,90)$,\allowbreak\ $(165,88)$\\
13 & 3 & $0$ & $(1,1902)$,\allowbreak\ $(44,2052)$,\allowbreak\ $(148,701)$,\allowbreak\ $(158,738)$,\allowbreak\ $(170,381)$,\allowbreak\ $(213,1376)$,\allowbreak\ $(317,1546)$,\allowbreak\ $(327,569)$,\allowbreak\ $(339,1057)$,\allowbreak\ $(382,700)$,\allowbreak\ $(486,194)$,\allowbreak\ $(496,400)$,\allowbreak\ $(508,1733)$,\allowbreak\ $(551,24)$,\allowbreak\ $(655,1039)$,\allowbreak\ $(665,231)$,\allowbreak\ $(677,212)$,\allowbreak\ $(720,1545)$,\allowbreak\ $(824,1884)$,\allowbreak\ $(834,62)$,\allowbreak\ $(846,888)$,\allowbreak\ $(889,869)$,\allowbreak\ $(993,532)$,\allowbreak\ $(1003,2090)$,\allowbreak\ $(1015,1564)$,\allowbreak\ $(1058,193)$,\allowbreak\ $(1162,1377)$,\allowbreak\ $(1172,1921)$,\allowbreak\ $(1184,43)$,\allowbreak\ $(1227,1714)$,\allowbreak\ $(1331,25)$,\allowbreak\ $(1341,1752)$,\allowbreak\ $(1353,719)$,\allowbreak\ $(1396,1038)$,\allowbreak\ $(1500,870)$,\allowbreak\ $(1510,1583)$,\allowbreak\ $(1522,1395)$,\allowbreak\ $(1565,362)$,\allowbreak\ $(1669,1715)$,\allowbreak\ $(1679,1414)$,\allowbreak\ $(1691,2071)$,\allowbreak\ $(1734,1883)$,\allowbreak\ $(1838,363)$,\allowbreak\ $(1848,1245)$,\allowbreak\ $(1860,550)$,\allowbreak\ $(1903,1207)$,\allowbreak\ $(2007,1208)$,\allowbreak\ $(2017,1076)$,\allowbreak\ $(2029,1226)$,\allowbreak\ $(2072,531)$,\allowbreak\ $(2176,2053)$,\allowbreak\ $(2186,907)$\\
\addlinespace[3pt]
\end{longtable}
\endgroup
\end{document}
